{\bfseries Caminata cuántica discreta con moneda de Fourier.
Considere la caminata cuántica discreta sobre la línea infinita de enteros. 
La moneda utilizada es la matriz definida por: 


\begin{equation*}
C = 
\begin{pmatrix} \frac{\sqrt{3}}{2} & \frac{1}{2} \\ \frac{1}{2} & -\frac{\sqrt{3}}{2} \end{pmatrix} 
\end{equation*}

El operador de desplazamiento es el habitual:
\begin{equation*}
S = \sum_{x \in \mathbb{Z}} (\ket{x+1}\bra{x} \otimes \ket{0}\bra{0} + \ket{x-1}\bra{x} \otimes \ket{1}\bra{1}) 
\end{equation*}

El estado inicial es $\ket{\psi(0)} = \ket{0}_p \otimes \ket{0}_c$ (partícula localizada en el
origen con moneda en $\ket{0}$). }

\begin{itemize}
    \item {\bfseries Aplique el paso evolutivo $U = S(I_p \otimes C)$ sucesivamente para obtener los
        estados $\ket{\psi(1)}, \ket{\psi(2)}, \ket{\psi(3)}, \ket{\psi(4)}, \ket{\psi(5)}$. }

    \item  {\bfseries Para cada paso, escriba explícitamente las amplitudes de probabilidad de las
            posiciones involucradas y calcule las probabilidades de encontrar la partícula en cada
        posición.}
    \item {\bfseries Comente si observa alguna diferencia cualitativa con la caminata de Hadamard
        (por ejemplo, simetría, velocidad de propagación, etc.).}
\end{itemize}

\vspace{.3cm}

\textbf{Paso 1:}
Aplicando el operador $U$ al estado inicial:
\begin{equation*}
\ket{\psi(1)} = S(I_p \otimes C) \ket{0,0} = S \left( \frac{\sqrt{3}}{2}\ket{0,0} + \frac{1}{2}\ket{0,1} \right) = \frac{\sqrt{3}}{2}\ket{1,0} + \frac{1}{2}\ket{-1,1}
\end{equation*}

Como nota importante, podemos pensar en la acción del desplazamiento como: 

\begin{align*}
    \ket{x,0} &\Rightarrow \ket{x+1,0} \\ 
    \ket{x,1} &\Rightarrow \ket{x-1,1} \\ 
\end{align*}

Y la acción de la moneda solo es multiplicación de matrices y multiplicación escalar, voy a omitir
el desarrollo para evitar redundancias.

Probabilidades:
\begin{align*}
P(1) &= \left| \frac{\sqrt{3}}{2} \right|^2 = \frac{3}{4} \\
P(-1) &= \left| \frac{1}{2} \right|^2 = \frac{1}{4}
\end{align*}

\textbf{Paso 2:}
\begin{equation*}
\ket{\psi(2)} = \frac{3}{4}\ket{2,0} + \frac{\sqrt{3}}{4}\ket{0,1} + \frac{1}{4}\ket{0,0} - \frac{\sqrt{3}}{4}\ket{-2,1}
\end{equation*}

Probabilidades:
\begin{align*}
P(2) &= \left| \frac{3}{4} \right|^2 = \frac{9}{16} \\
P(0) &= \left| \frac{1}{4} \right|^2 + \left| \frac{\sqrt{3}}{4} \right|^2 = \frac{4}{16} = \frac{1}{4} \\
P(-2) &= \left| -\frac{\sqrt{3}}{4} \right|^2 = \frac{3}{16}
\end{align*}

\textbf{Paso 3:}
\begin{equation*}
\ket{\psi(3)} = \frac{3\sqrt{3}}{8}\ket{3,0} + \frac{2\sqrt{3}}{8}\ket{1,0} + \frac{3}{8}\ket{1,1} - \frac{\sqrt{3}}{8}\ket{-1,0} - \frac{2}{8}\ket{-1,1} + \frac{3}{8}\ket{-3,1}
\end{equation*}

Probabilidades:
\begin{align*}
P(3) &= \frac{27}{64} \\
P(1) &= \frac{12}{64} + \frac{9}{64} = \frac{21}{64} \\
P(-1) &= \frac{3}{64} + \frac{4}{64} = \frac{7}{64} \\
P(-3) &= \frac{9}{64}
\end{align*}

\textbf{Paso 4:}
\begin{equation*}
\ket{\psi(4)} = \frac{9}{16}\ket{4,0} + \frac{9}{16}\ket{2,0} + \frac{3\sqrt{3}}{16}\ket{2,1} - \frac{5}{16}\ket{0,0} - \frac{\sqrt{3}}{16}\ket{0,1} + \frac{3}{16}\ket{-2,0} + \frac{\sqrt{3}}{16}\ket{-2,1} - \frac{3\sqrt{3}}{16}\ket{-4,1}
\end{equation*}

Probabilidades:
\begin{align*}
P(4) &= \frac{81}{256} \\
P(2) &= \frac{81}{256} + \frac{27}{256} = \frac{108}{256} = \frac{27}{64} \\
P(0) &= \frac{25}{256} + \frac{3}{256} = \frac{28}{256} = \frac{7}{64} \\
P(-2) &= \frac{9}{256} + \frac{3}{256} = \frac{12}{256} = \frac{3}{64} \\
P(-4) &= \frac{27}{256}
\end{align*}

\textbf{Paso 5:}
\begin{equation*}
\ket{\psi(5)} = \frac{9\sqrt{3}}{32}\ket{5,0} + \frac{12\sqrt{3}}{32}\ket{3,0} + \frac{9}{32}\ket{3,1} - \frac{6\sqrt{3}}{32}\ket{1,0} + \frac{4\sqrt{3}}{32}\ket{-1,0} - \frac{2}{32}\ket{-1,1} - \frac{3\sqrt{3}}{32}\ket{-3,0} + \frac{9}{32}\ket{-5,1}
\end{equation*}

Probabilidades:
\begin{align*}
P(5) &= \frac{243}{1024} \\
P(3) &= \frac{432}{1024} + \frac{81}{1024} = \frac{513}{1024} \\
P(1) &= \frac{108}{1024} = \frac{27}{256} \\
P(-1) &= \frac{48}{1024} + \frac{4}{1024} = \frac{52}{1024} = \frac{13}{256} \\
P(-3) &= \frac{27}{1024} \\
P(-5) &= \frac{81}{1024}
\end{align*}

\vspace{0.5cm}
\textbf{Diferencias cualitativas con la caminata de Hadamard:}
\begin{itemize}
    \item \textbf{Asimetría y sesgo:} La matriz dada introduce probabilidades desiguales ($3/4$ de
        probabilidad de asociarse a $\ket{0}_c$ y $1/4$ para $\ket{1}_c$). Como $\ket{0}_c$ desplaza
        la partícula hacia la derecha, el sistema presenta una asimetria hacia las posiciones positivas.
    \item \textbf{Propagación dirigida:} En $t=5$, cerca del $73.8\%$ de la distribución de probabilidad total se concentra en las posiciones $x=3$ y $x=5$. A diferencia de la moneda de Hadamard que distribuye la probabilidad más equitativamente (aunque con asimetría ), esta moneda reduce dramáticamente la interferencia destructiva en la dirección positiva, acelerando la localización de la partícula en ese frente de onda.
\end{itemize}

\vspace{.5cm}
